\section{Checklist：告诉 Agent ``什么叫好''}

\subsection{你在整个过程中唯一的工作}

\begin{importantbox}{你唯一要做的事}
给 agent 一份简单的 checklist，定义``什么叫好''。这是你在\textbf{整个 autoresearch 过程中唯一需要参与的环节}。

Checklist 由简单的是/否问题组成，每个问题检查输出的一个具体方面。通过或失败，就这么简单。Agent 用这份 checklist 对每次输出打分，分数变化告诉它改动是在帮忙还是在帮倒忙。
\end{importantbox}

\subsection{为什么用是/否而不用打分}

作者用``老师批卷子''来说明为什么要用二元判断而非主观评分。

不是``给写作质量打个 1--10 分''（模糊、每次结果不一样），而是每一项都清清楚楚是或否：

\begin{itemize}
    \item 学生有没有写论点陈述（thesis statement）？是或否。
    \item 每处引用都注明出处了吗？是或否。
    \item 篇幅在 5 页以内吗？是或否。
\end{itemize}

用这份 checklist 批 100 份卷子，每次结果都一致。

\begin{knowledgebox}{一致性是关键}
主观评分（如 1--10 分）在不同时间、不同上下文下会产生不同结果，导致评估信号噪声大。是/否二元判断消除了这种不确定性，让 agent 能可靠地判断改动的效果方向。
\end{knowledgebox}

\subsection{落地页文案 Checklist 示例}

以下是作者给出的落地页文案 skill 的具体 checklist 问题：

\begin{enumerate}
    \item ``标题有没有包含具体数字或结果？''——捕捉``Grow Your Business''这类模糊标题。
    \item ``文案中有没有 revolutionary、synergy、cutting-edge、next-level 这类流行词？''——捕捉空洞的 buzzword 堆砌。
    \item ``CTA 是否使用了具体的动词短语？''——捕捉``Learn More''或``Click Here''这类弱 CTA。
    \item ``第一句话有没有点出一个具体的痛点？''——捕捉``In today's fast-paced world...''这类泛泛的开场。
    \item ``总文案字数是否在 150 词以内？''——捕捉臃肿的页面。
\end{enumerate}

\subsection{Checklist 的生成与数量}

这些 checklist 不需要你自己费劲想。启动 autoresearch 时，agent 会引导你完成整个过程：

\begin{itemize}
    \item 问你``什么叫好''
    \item 帮你把模糊的感觉（vibes）变成具体的是/否问题
    \item 如果你有现成的 style guide，还能从中提取标准
\end{itemize}

\begin{warningbox}{Checklist 数量的陷阱}
\textbf{3--6 个问题是最佳数量。}超过这个范围，skill 会开始``应付 checklist''——就像学生背答案而不理解题目。这是一种 overfitting 现象：skill 在 checklist 上得分越来越高，但实际输出质量反而下降。
\end{warningbox}

\subsection{本章小结}

评分标准是 autoresearch 中唯一需要人类定义的部分。使用是/否二元判断（而非主观评分）确保评估的一致性和可复现性。Checklist 应控制在 3--6 个问题，避免 overfitting。Agent 会引导你生成 checklist，你不需要从零设计。

